% Options for packages loaded elsewhere
\PassOptionsToPackage{unicode}{hyperref}
\PassOptionsToPackage{hyphens}{url}
\documentclass[
]{article}
\usepackage{xcolor}
\usepackage[margin=1in]{geometry}
\usepackage{amsmath,amssymb}
\setcounter{secnumdepth}{-\maxdimen} % remove section numbering
\usepackage{iftex}
\ifPDFTeX
  \usepackage[T1]{fontenc}
  \usepackage[utf8]{inputenc}
  \usepackage{textcomp} % provide euro and other symbols
\else % if luatex or xetex
  \usepackage{unicode-math} % this also loads fontspec
  \defaultfontfeatures{Scale=MatchLowercase}
  \defaultfontfeatures[\rmfamily]{Ligatures=TeX,Scale=1}
\fi
\usepackage{lmodern}
\ifPDFTeX\else
  % xetex/luatex font selection
\fi
% Use upquote if available, for straight quotes in verbatim environments
\IfFileExists{upquote.sty}{\usepackage{upquote}}{}
\IfFileExists{microtype.sty}{% use microtype if available
  \usepackage[]{microtype}
  \UseMicrotypeSet[protrusion]{basicmath} % disable protrusion for tt fonts
}{}
\makeatletter
\@ifundefined{KOMAClassName}{% if non-KOMA class
  \IfFileExists{parskip.sty}{%
    \usepackage{parskip}
  }{% else
    \setlength{\parindent}{0pt}
    \setlength{\parskip}{6pt plus 2pt minus 1pt}}
}{% if KOMA class
  \KOMAoptions{parskip=half}}
\makeatother
\usepackage{longtable,booktabs,array}
\newcounter{none} % for unnumbered tables
\usepackage{calc} % for calculating minipage widths
% Correct order of tables after \paragraph or \subparagraph
\usepackage{etoolbox}
\makeatletter
\patchcmd\longtable{\par}{\if@noskipsec\mbox{}\fi\par}{}{}
\makeatother
% Allow footnotes in longtable head/foot
\IfFileExists{footnotehyper.sty}{\usepackage{footnotehyper}}{\usepackage{footnote}}
\makesavenoteenv{longtable}
\setlength{\emergencystretch}{3em} % prevent overfull lines
\providecommand{\tightlist}{%
  \setlength{\itemsep}{0pt}\setlength{\parskip}{0pt}}
\usepackage{bookmark}
\IfFileExists{xurl.sty}{\usepackage{xurl}}{} % add URL line breaks if available
\urlstyle{same}
\hypersetup{
  hidelinks,
  pdfcreator={LaTeX via pandoc}}

\author{}
\date{}

\begin{document}

\section{\texorpdfstring{A pre-registered, falsifiable
measurement-methodology bridge from \texttt{styxx} primitives to EU AI
Act Article 15 / Annex III requirements
(v0.1)}{A pre-registered, falsifiable measurement-methodology bridge from styxx primitives to EU AI Act Article 15 / Annex III requirements (v0.1)}}\label{a-pre-registered-falsifiable-measurement-methodology-bridge-from-styxx-primitives-to-eu-ai-act-article-15-annex-iii-requirements-v0.1}

\textbf{Author:} Alexander Rodabaugh (Fathom Lab) \textbf{Date:}
2026-05-28 \textbf{Substrate:} styxx 7.7.10 in-development; companion to
\texttt{papers/PAPER\_recursive\_discipline\_2026\_05\_27.md} (v7);
module \texttt{styxx.compliance.eu\_ai\_act} ships in this same commit
\textbf{Status:} v0.1 minimum-viable mapping. Published BEFORE the EU AI
Act high-risk system enforcement deadline of 2 August 2026 to give
regulated operators an evaluation runway. Not legal advice. Independent
conformity review required for any production deployment.

\begin{center}\rule{0.5\linewidth}{0.5pt}\end{center}

\subsection{Abstract}\label{abstract}

The EU AI Act high-risk obligations enter enforcement on 2 August 2026
with penalties up to €15M or 3\% of global annual turnover. Article 15
mandates that high-risk AI systems achieve appropriate levels of
accuracy, robustness, and cybersecurity, that accuracy metrics be
declared in the instructions of use, and --- under paragraph 2 --- that
\emph{``the Commission shall, in cooperation with relevant
stakeholders\ldots{} encourage the development of benchmarks and
measurement methodologies''}. The Commission's invitation is open. No
competing AI observability or evaluation product currently publishes a
structured Article 15 mapping. This paper introduces
\texttt{styxx.compliance.eu\_ai\_act}, the first open-source,
pre-registration-disciplined measurement-methodology bridge mapping a
deployable cognitive-observability primitive set to specific Article 15
sub-paragraphs, with calibrated metrics, explicit construct-ceiling
disclosures, commit-level reproducibility receipts, and pre-stated
falsification criteria. The v0.1 mapping covers four Article 15 clauses
with five styxx primitives; it explicitly enumerates seven uncovered EU
AI Act requirements (Article 9, 10, 12, 13, 14, 15.4 bias, 15
cybersecurity) and points each at a non-styxx alternative tool or
methodology. The boundary statement is at least as long as the coverage
statement by design (kill-gate A3). Five pre-registered kill-gates
define what success and failure look like for the v0.1 release. This
document is a measurement methodology, not legal advice; it is the kind
of artifact Article 15 paragraph 2 explicitly invites stakeholders to
develop.

\begin{center}\rule{0.5\linewidth}{0.5pt}\end{center}

\subsection{1. Background}\label{background}

\subsubsection{1.1 The August 2, 2026
deadline}\label{the-august-2-2026-deadline}

The EU AI Act (Regulation (EU) 2024/1689) enters its high-risk system
enforcement phase on 2 August 2026. Annex III enumerates high-risk
categories including biometrics, critical infrastructure, education,
employment, essential services, law enforcement, migration, and
democratic process administration. Providers of high-risk systems must
demonstrate conformity through Article 9 risk management, Article 10
data governance, Article 11 technical documentation, Article 12 logging,
Article 13 transparency, Article 14 human oversight, and Article 15
accuracy/robustness/cybersecurity. Most Annex III systems may be
self-assessed where harmonised standards apply; biometric, critical
infrastructure, and law-enforcement systems require third-party
conformity assessment by a notified body. Penalties for non-compliance
reach €15,000,000 or 3\% of global annual turnover, whichever is higher.

\subsubsection{1.2 Article 15 in detail}\label{article-15-in-detail}

Article 15 (``Accuracy, robustness and cybersecurity'') has four
substantive paragraphs:

\begin{itemize}
\tightlist
\item
  \textbf{15.1}: high-risk AI systems shall be \emph{designed and
  developed} to achieve appropriate levels of accuracy, robustness, and
  cybersecurity \emph{throughout their lifecycle}.
\item
  \textbf{15.1(a)} (instructions of use): accuracy levels and
  \emph{relevant accuracy metrics} shall be declared in the accompanying
  instructions of use.
\item
  \textbf{15.2}: \emph{``the Commission shall, in cooperation with
  relevant stakeholders and organisations such as metrology and
  benchmarking authorities, encourage, as appropriate, the development
  of benchmarks and measurement methodologies''}. This is the open
  invitation that this paper responds to.
\item
  \textbf{15.3}: robustness shall be as high as possible regarding
  errors, faults, or inconsistencies, achievable through \emph{technical
  redundancy solutions} including backup or fail-safe plans.
\item
  \textbf{15.4}: high-risk AI systems that continue to learn after
  deployment shall be designed to \emph{eliminate or reduce as far as
  possible} the risk of biased outputs influencing future inputs
  (feedback loops).
\end{itemize}

\subsubsection{1.3 What this paper IS, and is
NOT}\label{what-this-paper-is-and-is-not}

This paper IS: - A measurement methodology: it names specific styxx
primitives, calibrated metrics, construct ceilings, and reproducibility
receipts that produce evidence relevant to specific Article 15 clauses.
- Pre-registered: every claim cites a commit, every primitive's failure
mode is disclosed, every uncovered requirement is enumerated with an
alternative tool reference. - Open-source and falsifiable: the
\texttt{styxx.compliance.eu\_ai\_act} module is MIT-licensed Python; the
tests verify the structural integrity of the mapping; the paper's claims
can be re-checked at any future git commit.

This paper IS NOT: - Legal advice. Conformity assessment requires legal
expertise this paper does not provide. - Sufficient on its own for EU AI
Act conformity. It addresses Article 15 sub-paragraphs only and
explicitly disclaims coverage of Articles 9, 10, 12, 13, 14. - A claim
that styxx primitives have been validated by a notified body. They have
been validated against the styxx project's own benchmarks; third-party
validation remains future work. - An endorsement by the European
Commission, AISI, METR, Apollo Research, or any standardization body. It
is a unilateral stakeholder contribution.

\begin{center}\rule{0.5\linewidth}{0.5pt}\end{center}

\subsection{2. Methodology}\label{methodology}

\subsubsection{2.1 The bridge primitives}\label{the-bridge-primitives}

\texttt{styxx.compliance.eu\_ai\_act} exposes four objects:

\begin{enumerate}
\def\labelenumi{\arabic{enumi}.}
\tightlist
\item
  \textbf{\texttt{ARTICLE\_15\_REGISTRY:\ dict{[}str,\ ComplianceMap{]}}}
  --- the registry of mapped clauses. v0.1 keys are
  \texttt{"Article\ 15.1"}, \texttt{"Article\ 15.1(a)"},
  \texttt{"Article\ 15.3"}, \texttt{"Article\ 15.4"}.
\item
  \textbf{\texttt{ComplianceMap}} --- dataclass with \texttt{clause},
  \texttt{requirement\_text}, \texttt{styxx\_primitives} (tuple of
  \texttt{PrimitiveCoverage}), and \texttt{notes}. Each \texttt{notes}
  field includes a \emph{``not legal advice''} disclaimer.
\item
  \textbf{\texttt{PrimitiveCoverage}} --- per-primitive coverage entry
  with \texttt{primitive} (public API symbol),
  \texttt{calibrated\_metric} (published AUC/accuracy with context),
  \texttt{construct\_ceiling} (honest failure mode),
  \texttt{receipt\_commit} (styxx git SHA that produced the metric), and
  \texttt{receipt\_doc} (path to the validating FINDING).
\item
  \textbf{\texttt{UNCOVERED\_REQUIREMENTS:\ tuple{[}UncoveredRequirement,\ ...{]}}}
  --- the boundary statement: every EU AI Act requirement styxx does NOT
  cover, with reason and alternative tool reference.
\end{enumerate}

Helper functions: \texttt{cite(article:\ str)},
\texttt{coverage\_table()}, \texttt{uncovered\_requirements()}.

\subsubsection{2.2 Five pre-registered
kill-gates}\label{five-pre-registered-kill-gates}

These were pre-stated in the strategic landscape document
(\texttt{.styxx/STRATEGIC\_LANDSCAPE\_2026\_05\_28.md}, written before
this module existed) and are enforced by the test suite
(\texttt{tests/test\_compliance\_eu\_ai\_act.py}):

\begin{itemize}
\tightlist
\item
  \textbf{A1 (validity).} Every clause key in
  \texttt{ARTICLE\_15\_REGISTRY} must cite a specific Article
  sub-paragraph, not generic compliance language. Enforced by regex
  \texttt{\^{}Article\ \textbackslash{}d+(\textbackslash{}.\textbackslash{}d+)?(\textbackslash{}({[}a-z{]}\textbackslash{}))?\$}.
\item
  \textbf{A2 (falsifiability).} Every
  \texttt{PrimitiveCoverage.construct\_ceiling} must be non-empty and at
  least 50 characters. Hidden caveats are not permitted; the failure
  mode goes in the field.
\item
  \textbf{A3 (boundary explicitness).}
  \texttt{len(uncovered\_requirements())\ \textgreater{}=\ len(coverage\_table())}.
  The list of what styxx does NOT cover must be at least as long as the
  list of what it does. v0.1 ships 7 uncovered vs 4 covered.
\item
  \textbf{A4 (timeline).} Publish before July 1, 2026 to provide a
  \textasciitilde30-day evaluation runway before the August 2
  enforcement deadline. This paper is committed 2026-05-28.
\item
  \textbf{A5 (citation strategy).} ≥1 independent citation (academic,
  regulator, enterprise) within 6 months of publication. If 0 by
  2027-02-01, the methodology did not achieve uptake and is reassessed.
  The strategic landscape doc references this measurable kill-gate
  explicitly.
\end{itemize}

A2 and A3 are enforced at test time. A1 is enforced at test time. A4 and
A5 are enforced by the project timeline.

\subsubsection{2.3 What the mapping does NOT
attempt}\label{what-the-mapping-does-not-attempt}

\begin{itemize}
\tightlist
\item
  It does not score the \textbf{adequacy} of styxx primitives for any
  specific high-risk Annex III system. That assessment is
  system-specific and operator-specific.
\item
  It does not enumerate ALL Article 15 sub-paragraphs. v0.1 covers four.
  v0.2 may extend.
\item
  It does not write conformity declarations. Operators write
  declarations; styxx supplies the measurement evidence.
\item
  It does not claim to satisfy harmonised standards (none for AI agent
  honesty exist as of 2026-05-28).
\end{itemize}

\begin{center}\rule{0.5\linewidth}{0.5pt}\end{center}

\subsection{3. The coverage table}\label{the-coverage-table}

Each row below corresponds to one entry in
\texttt{ARTICLE\_15\_REGISTRY}. Every primitive in the right column
ships in styxx 7.7.10; every commit hash is reachable at
\texttt{fathom-lab/styxx@main}.

\subsubsection{3.1 Article 15.1 --- appropriate accuracy, robustness,
cybersecurity throughout
lifecycle}\label{article-15.1-appropriate-accuracy-robustness-cybersecurity-throughout-lifecycle}

{\def\LTcaptype{none} % do not increment counter
\begin{longtable}[]{@{}
  >{\raggedright\arraybackslash}p{(\linewidth - 4\tabcolsep) * \real{0.3333}}
  >{\raggedright\arraybackslash}p{(\linewidth - 4\tabcolsep) * \real{0.3333}}
  >{\raggedright\arraybackslash}p{(\linewidth - 4\tabcolsep) * \real{0.3333}}@{}}
\toprule\noalign{}
\begin{minipage}[b]{\linewidth}\raggedright
primitive
\end{minipage} & \begin{minipage}[b]{\linewidth}\raggedright
calibrated metric
\end{minipage} & \begin{minipage}[b]{\linewidth}\raggedright
construct ceiling
\end{minipage} \\
\midrule\noalign{}
\endhead
\bottomrule\noalign{}
\endlastfoot
\texttt{styxx.cognometric\_card} & HaluEval-QA AUC 0.998 (mean over 150
items, seeds 31/47/83); XSTest refusal AUC 0.976; BFCL v3 tool-drift AUC
0.943 & Text-only register space has construct ceilings: instruments
measure register, not always content. Sycophancy false-positive 0.30 on
restrained-tech responses (gpt-3.5: 0.60). Logprob-validity works for
refusal but fails for hallucination (confident confabulation). \\
\texttt{styxx.gauntlet\ +\ styxx.preflight} & v3 detection bars
(D1+D2+D3+D4) with regression-tested oracles. 18 pre-registered
baselines tested; 17 FAILed bars before Baseline-019 PASSed. & Bars
catch confounds the project pre-stated; do not catch unknown-unknown
confound classes. Each bar revision documents one prior missed
confound. \\
\texttt{styxx.agent\_audit} & 13/13 PASS modal pre-stated outcome (L5
commit \texttt{3c24b5e}); L7 uncurated extension caught off-by-one count
drift in companion paper (2 FAILs as pre-disclosed). &
First-occurrence-only by default --- caught initial drift but missed
propagation to 4 additional places. Richer multi-occurrence checker is
methodology future work. \\
\end{longtable}
}

\subsubsection{3.2 Article 15.1(a) --- accuracy metrics in instructions
of use}\label{article-15.1a-accuracy-metrics-in-instructions-of-use}

{\def\LTcaptype{none} % do not increment counter
\begin{longtable}[]{@{}
  >{\raggedright\arraybackslash}p{(\linewidth - 4\tabcolsep) * \real{0.3333}}
  >{\raggedright\arraybackslash}p{(\linewidth - 4\tabcolsep) * \real{0.3333}}
  >{\raggedright\arraybackslash}p{(\linewidth - 4\tabcolsep) * \real{0.3333}}@{}}
\toprule\noalign{}
\begin{minipage}[b]{\linewidth}\raggedright
primitive
\end{minipage} & \begin{minipage}[b]{\linewidth}\raggedright
calibrated metric
\end{minipage} & \begin{minipage}[b]{\linewidth}\raggedright
construct ceiling
\end{minipage} \\
\midrule\noalign{}
\endhead
\bottomrule\noalign{}
\endlastfoot
\texttt{styxx.cognometric\_card} & per-step cognometric readout (drift,
confabulation, refusal, sycophancy, phase transition, low trust,
incoherence) & construct ceilings apply per-axis; see §3.1 \\
\texttt{styxx.critique\_detector(model=\textquotesingle{}gpt-4o-mini\textquotesingle{})}
& Baseline-019 PASSes gauntlet v3 bars at AUC 0.95, pre-stated 28\%
probability landed cleanly & Mechanism is ``out-of-context critique'',
NOT within-model generation-vs-critique asymmetry. True within-model
asymmetry: 5.88\% on dark-core / 17.00\% on TruthfulQA (v3 measurement).
In-council bias: default backend gpt-4o-mini was in the original
3-vendor council. \\
\texttt{styxx.gauntlet\ +\ styxx.preflight} & see §3.1 & see §3.1 \\
\end{longtable}
}

These three together produce the \emph{``relevant accuracy metrics''} an
operator can declare in instructions of use, with the honest
construct-ceiling disclosure that Article 15.1(a) does not explicitly
mandate but that the recursive-discipline thesis argues should be
present in any defensible declaration.

\subsubsection{3.3 Article 15.3 --- robustness via technical redundancy
/ fail-safe
plans}\label{article-15.3-robustness-via-technical-redundancy-fail-safe-plans}

{\def\LTcaptype{none} % do not increment counter
\begin{longtable}[]{@{}
  >{\raggedright\arraybackslash}p{(\linewidth - 4\tabcolsep) * \real{0.3333}}
  >{\raggedright\arraybackslash}p{(\linewidth - 4\tabcolsep) * \real{0.3333}}
  >{\raggedright\arraybackslash}p{(\linewidth - 4\tabcolsep) * \real{0.3333}}@{}}
\toprule\noalign{}
\begin{minipage}[b]{\linewidth}\raggedright
primitive
\end{minipage} & \begin{minipage}[b]{\linewidth}\raggedright
calibrated metric
\end{minipage} & \begin{minipage}[b]{\linewidth}\raggedright
construct ceiling
\end{minipage} \\
\midrule\noalign{}
\endhead
\bottomrule\noalign{}
\endlastfoot
\texttt{styxx.recover\_posture} & cognitive-integrity persistence
primitive for agents crossing context-compaction boundaries & v1 reads
what chart.jsonl persists; does not include cogn\_audit scores. Not
validated on third-party agent platforms. \\
\texttt{styxx.agent\_audit} & substrate-grounded session-output verifier
& see §3.1 \\
\end{longtable}
}

Article 15.3 envisions backup/fail-safe plans against runtime errors.
styxx's \texttt{recover\_posture} is a cognition-side primitive: it does
not replace a process-level fail-safe (an external watchdog, an
interrupt key, a rollback procedure) but it provides verifiable evidence
of cognitive-integrity continuity across the most common modern failure
mode (context-window compaction in long-running agents).

\subsubsection{3.4 Article 15.4 --- bias amplification / feedback-loop
mitigation}\label{article-15.4-bias-amplification-feedback-loop-mitigation}

{\def\LTcaptype{none} % do not increment counter
\begin{longtable}[]{@{}lll@{}}
\toprule\noalign{}
primitive & calibrated metric & construct ceiling \\
\midrule\noalign{}
\endhead
\bottomrule\noalign{}
\endlastfoot
\emph{(none --- honest empty coverage)} & --- & --- \\
\end{longtable}
}

v0.1 has \textbf{NO} styxx-side coverage for Article 15.4. This empty
cell is intentional and is itself a discipline statement: operators must
consult separate tools (fairness libraries, drift monitors).
\texttt{styxx.cognometric\_card} provides per-step register signals that
\emph{may} surface drift indirectly but is not a substitute. See §4 for
the boundary alternative reference.

\begin{center}\rule{0.5\linewidth}{0.5pt}\end{center}

\subsection{4. What styxx does NOT cover (the boundary
statement)}\label{what-styxx-does-not-cover-the-boundary-statement}

Per kill-gate A3, this section is intentionally longer than §3. Seven EU
AI Act clauses are explicitly enumerated as outside the v0.1 scope, with
alternative tools pointed at:

\subsubsection{4.1 Article 15.4 --- bias
amplification}\label{article-15.4-bias-amplification}

styxx primitives operate on per-step agent cognition signals; they do
not measure population-level disparate impact, demographic outcome
equalization, or training-data bias amplification across protected
classes. \textbf{Alternative}: Fairness audit libraries (Fairlearn,
AIF360) plus domain-specific impact assessment; legal review.

\subsubsection{4.2 Article 15 ---
cybersecurity}\label{article-15-cybersecurity}

styxx instruments observe agent cognition, not network, host, or
supply-chain security. Prompt-injection robustness is \emph{adjacent}
(refusal-related instruments may flag some injection attempts) but not
the primary scope. \textbf{Alternative}: dedicated LLM security tools
(Lakera Guard, Prompt-Armor, OWASP LLM Top 10 mitigations); standard
application-security audit; penetration testing.

\subsubsection{4.3 Article 9 --- risk
management}\label{article-9-risk-management}

styxx provides MEASUREMENT evidence; it does not implement the
risk-management lifecycle (identification, estimation, evaluation,
mitigation, residual-risk acceptance) prescribed by Article 9.
\textbf{Alternative}: ISO/IEC 23894:2023 or NIST AI RMF 1.0 frameworks;
QMS-style organizational processes.

\subsubsection{4.4 Article 10 --- data
governance}\label{article-10-data-governance}

styxx does not score training-data quality, provenance, labeling
consistency, or representativeness. Cognometric instruments operate
post-training on agent outputs. \textbf{Alternative}: data documentation
frameworks (Datasheets for Datasets, Data Statements); training-data
audit tooling.

\subsubsection{4.5 Article 12 ---
record-keeping}\label{article-12-record-keeping}

styxx writes per-step cognometric vitals but does not, by itself,
satisfy Article 12's logging-traceability requirements end-to-end (event
capture, retention, integrity). \textbf{Alternative}: production
observability platforms (OpenTelemetry, Langfuse, Arize, Datadog LLM
Observability) for trace-level logging alongside styxx for
cognition-level scoring.

\subsubsection{4.6 Article 13 --- transparency to
deployers}\label{article-13-transparency-to-deployers}

styxx instruments produce evidence; they do not generate the
deployer-facing instructions of use, capability statements, or
output-explanation interfaces that Article 13 mandates.
\textbf{Alternative}: model card frameworks; capability-statement
templates; deployer documentation processes.

\subsubsection{4.7 Article 14 --- human
oversight}\label{article-14-human-oversight}

styxx is a measurement layer; it does not implement the
human-in-the-loop interfaces, stop-button affordances, or
operator-control surfaces Article 14 requires. \textbf{Alternative}:
agent-platform-level oversight UI; interrupt and rollback primitives in
the agent runtime; operator training programs.

\begin{center}\rule{0.5\linewidth}{0.5pt}\end{center}

\subsection{5. Pre-registered falsification criteria for THIS
paper}\label{pre-registered-falsification-criteria-for-this-paper}

The recursive-discipline methodology (companion paper,
\texttt{papers/PAPER\_recursive\_discipline\_2026\_05\_27.md} v7) argues
that papers should pre-register their own falsification criteria. This
paper does:

\begin{itemize}
\tightlist
\item
  \textbf{F1} (registry validity): if any clause in
  \texttt{ARTICLE\_15\_REGISTRY} cites an Article paragraph that does
  not exist in the EU AI Act final consolidated text, this paper is
  falsified. \textbf{Verification}: regex-checked at test time + manual
  citation review.
\item
  \textbf{F2} (calibrated-metric reproducibility): if any cited
  AUC/accuracy number cannot be reproduced from the cited commit hash
  within ±0.01, that PrimitiveCoverage entry is falsified.
  \textbf{Verification}: re-run
  \texttt{submissions/\_results/leaderboard.json} reproduction scripts.
\item
  \textbf{F3} (construct-ceiling discipline): if any reviewer can
  produce a published failure mode of a styxx primitive that is NOT
  mentioned in any \texttt{PrimitiveCoverage.construct\_ceiling} field,
  that entry is falsified and must be updated. \textbf{Verification}:
  open issue acceptance protocol.
\item
  \textbf{F4} (boundary violation): if a styxx primitive is later shown
  to produce evidence for an Article in
  \texttt{UNCOVERED\_REQUIREMENTS}, the mapping was scoped too narrowly.
  \textbf{Verification}: re-evaluation by versioned issue.
\item
  \textbf{F5} (citation: A5): no independent citation by 2027-02-01
  means the methodology did not achieve uptake; reassess relevance and
  consider deprecation.
\end{itemize}

Every falsification path is observable and has a defined response.

\begin{center}\rule{0.5\linewidth}{0.5pt}\end{center}

\subsection{6. Reproducibility appendix}\label{reproducibility-appendix}

{\def\LTcaptype{none} % do not increment counter
\begin{longtable}[]{@{}
  >{\raggedright\arraybackslash}p{(\linewidth - 4\tabcolsep) * \real{0.3333}}
  >{\raggedright\arraybackslash}p{(\linewidth - 4\tabcolsep) * \real{0.3333}}
  >{\raggedright\arraybackslash}p{(\linewidth - 4\tabcolsep) * \real{0.3333}}@{}}
\toprule\noalign{}
\begin{minipage}[b]{\linewidth}\raggedright
artifact
\end{minipage} & \begin{minipage}[b]{\linewidth}\raggedright
commit
\end{minipage} & \begin{minipage}[b]{\linewidth}\raggedright
path
\end{minipage} \\
\midrule\noalign{}
\endhead
\bottomrule\noalign{}
\endlastfoot
this paper & this commit &
\texttt{papers/EU\_AI\_ACT\_COMPLIANCE\_2026.md} \\
compliance module package & this commit &
\texttt{styxx/compliance/\_\_init\_\_.py},
\texttt{styxx/compliance/eu\_ai\_act.py},
\texttt{styxx/compliance/\_legacy.py} (preserved v1.3.0 API) \\
compliance tests & this commit &
\texttt{tests/test\_compliance\_eu\_ai\_act.py} (15 tests, all enforce
A1--A3 kill-gates) \\
strategic landscape (pre-stated kill-gates) & private to operator &
\texttt{.styxx/STRATEGIC\_LANDSCAPE\_2026\_05\_28.md} \\
Baseline-019 result (Article 15.1(a) AUC 0.95 receipt) &
\texttt{17fdd97} &
\texttt{submissions/baseline\_019\_openai\_critique/submission.json} \\
Asymmetry-v3 result (mechanism correction receipt) & \texttt{ed663ca} &
\texttt{experiments/asymmetry\_v3\_cleanup\_2026\_05\_27/results.json} \\
L5 agent\_audit run (Article 15.1 evidence) & \texttt{3c24b5e} &
\texttt{experiments/agent\_claim\_audit\_2026\_05\_28/results.json} \\
L7 uncurated audit (boundary discipline receipt) & \texttt{cf14c83} &
\texttt{experiments/v6\_uncurated\_audit\_2026\_05\_28/results.json} \\
recursive-discipline paper v7 (companion) & \texttt{cf14c83} &
\texttt{papers/PAPER\_recursive\_discipline\_2026\_05\_27.md} \\
\end{longtable}
}

All commit hashes resolve at \texttt{fathom-lab/styxx@main}. The
\texttt{styxx.compliance.eu\_ai\_act} module's structural integrity is
verified by \texttt{tests/test\_compliance\_eu\_ai\_act.py} at every CI
run.

\begin{center}\rule{0.5\linewidth}{0.5pt}\end{center}

\subsection{7. Limitations + scope of
applicability}\label{limitations-scope-of-applicability}

\begin{enumerate}
\def\labelenumi{\arabic{enumi}.}
\tightlist
\item
  \textbf{v0.1 covers four Article 15 clauses only.} Articles 15.2
  (Commission methodology development), 15.5 (high-risk performance
  specifications), and Annex IV documentation requirements are out of
  scope for v0.1.
\item
  \textbf{No notified-body endorsement.} Self-assessment paths under
  Article 43 may accept this methodology as supporting evidence, but
  third-party conformity assessment (biometrics, critical
  infrastructure, law enforcement) requires notified-body review.
\item
  \textbf{Single-substrate validation.} styxx primitives have been
  validated against the styxx project's own benchmarks (HaluEval-QA,
  XSTest, BFCL v3, dark-core, TruthfulQA). External validation on
  customer deployments is future work.
\item
  \textbf{Open question on harmonised standards.} As of 2026-05-28 there
  is no harmonised standard for AI agent honesty or sycophancy
  measurement under the EU AI Act. CEN-CENELEC JTC 21 is the relevant
  standardization body; submission to JTC 21 requires institutional
  sponsorship and is operator territory, not styxx-project unilateral
  action.
\item
  \textbf{Cross-jurisdiction}: this methodology addresses EU AI Act
  Article 15 only. It does NOT address US (NIST AI RMF voluntary), UK
  (AISI guidance), or other jurisdictional regimes, though many
  requirements have functional analogues.
\end{enumerate}

\begin{center}\rule{0.5\linewidth}{0.5pt}\end{center}

\subsection{8. Conclusion}\label{conclusion}

\texttt{styxx.compliance.eu\_ai\_act} is, to the author's knowledge as
of 2026-05-28, the first open-source measurement-methodology bridge
mapping a deployable AI agent cognitive-observability primitive set to
specific EU AI Act Article 15 sub-paragraphs. v0.1 covers four clauses
with five primitives, enumerates seven uncovered requirements with
alternative tool references, ships under MIT license alongside the
underlying primitives, enforces structural-integrity kill-gates at test
time, and pre-registers five falsification criteria with defined
response paths. The bridge is published before the 2 August 2026
enforcement deadline to give regulated operators an evaluation runway.

The methodology does not satisfy EU AI Act conformity on its own; it is
one component, honestly bounded. Operators must conduct independent
legal review, apply harmonised standards where they exist, and consult
the alternative tools enumerated in §4 for the seven EU AI Act
requirements styxx does NOT cover.

What this paper is: the kind of stakeholder methodology contribution
that Article 15 paragraph 2 explicitly invites. What it is not: a
substitute for an operator's own conformity assessment.

\begin{center}\rule{0.5\linewidth}{0.5pt}\end{center}

\subsection{Acknowledgments}\label{acknowledgments}

Written 2026-05-28 alongside \texttt{styxx.compliance.eu\_ai\_act} v0.1
in a continuous session, with the cognitive support of Claude Opus 4.7
acting as in-session collaborator. The recursive-discipline methodology
of the companion paper
(\texttt{PAPER\_recursive\_discipline\_2026\_05\_27.md} v7) directly
informs this paper's pre-registration kill-gates, the construct-ceiling
discipline in §3, and the boundary-statement discipline in §4. Every
claim in this paper is reproducible at commit-level granularity from
\texttt{fathom-lab/styxx@main}.

This paper is offered to the European Commission, AISI, METR, Apollo
Research, FAR.AI, and CEN-CENELEC JTC 21 as an open-source stakeholder
contribution under Article 15 paragraph 2. Comments, falsifications, and
improvements are explicitly invited via the \texttt{fathom-lab/styxx}
GitHub issue tracker.

\end{document}
